\documentclass[11pt]{article}
\usepackage[margin=1in]{geometry}
\usepackage[T1]{fontenc}
\usepackage[utf8]{inputenc}
\usepackage{amsmath,amssymb}
\usepackage{booktabs,longtable}
\usepackage{enumitem}
\usepackage[hidelinks]{hyperref}

\title{Ratcheting Consequence Inquiry\\
Normative Project Specification}
\author{RCI Project}
\date{Version 0.5 --- 2026-08-22}

\begin{document}
\maketitle

\section*{Status, interpretation, and conformance}

This document is the normative semantic specification for Ratcheting Consequence
Inquiry (RCI) version 0.5. It preserves the v0.4 aggregate/history correction and is a
reconciled rewrite of the archived research
drafts listed in \path{docs/source-manifest.md}. Those drafts are provenance and
research context, not live instructions. Direct user instructions retain runtime
precedence; \path{PLAN.md} fixes approved engineering choices and delivery order;
the active Goal fixes the current completion boundary.

The key words \textbf{MUST}, \textbf{MUST NOT}, \textbf{SHOULD}, and
\textbf{MAY} are normative. A numbered requirement is conformant only when
\path{docs/requirements-matrix.md} points to checked evidence or explicitly
assigns it to a later milestone or open research question. Interfaces named for
later phases specify semantic extension seams; their mere presence does not
claim executable support.

\section{Source reconciliation}

\begin{longtable}{p{0.18\linewidth}p{0.30\linewidth}p{0.44\linewidth}}
\toprule
Source & Role in 0.5 & Reconciliation \\
\midrule
\endhead
RCI 0.1 & Original project and staged implementation draft &
Superseded as a live specification. Its claim/fact boundary, formal attack
discipline, phase separation, and deterministic reference-system intent are
retained. \\
RCI 0.2 & Question-conditioned architecture &
Its consequence-relative equivalence, stretch/squeeze discipline, typed
questions, non-explosion, lazy formalization, and guarded warrant are retained. \\
RCI 0.3 cognitive integration & Cognitive and memory research draft &
Recurrent probes, prediction seals, action--return contingency, reconstruction,
episodic/semantic separation, and retention are retained after removing duplicate
authorities and separating actual returns from interpretation. \\
0.3 Repair Delta & Corrective architecture note &
Support environments, route-specific open dependency, explicit attempt outcomes,
return routes, raw returns, derived working views, versioned ancestry, and
compression residue are adopted. Cyclic proof structures are admitted only as
inert data unless a later independently checked proof discipline is authorized;
the active policy remains cycle-free. \\
Consequence-subspace note & Compression and linear-binding note &
The exact/approximate and object/consequence-regeneration distinctions are
adopted. The linear theorem is corrected in Requirements RCI-051--RCI-055.
Compression execution remains Phase 3 work. \\
Retention/reacquisition note & Recovery and relearning note &
Present use, reconstruction, direct consequence evaluation, and reacquisition
are retained as independent relations. Reacquisition advantage is meaningful
only against a pinned baseline and cost ordering. Executable recovery policy is
G2/G3 work, not G1 acceptance. \\
Opaque controlled-memory environment & Future end-to-end benchmark note &
Its permanent epistemic distinctions and staged capability requirements are
retained. The environment is a G7 benchmark, not a core primitive or current
test. Predictive-state representation is a future native binding, and the
linear theorem applies only after a suitable representation has been
established. \\
AGENTS and Goal drafts & Governance proposals &
Archived for provenance. They do not override this document, the root
\path{AGENTS.md}, \path{PLAN.md}, or an active Goal. \\
RCI 0.4 repository-aligned semantic specification & Aggregate/history/state
correction and G3 target &
Adopted as the semantic correction governing G3. It distinguishes the
replay-complete aggregate fold from a binding-defined consequence-sufficient
history representation and preserves all sealed G1, G2A, and G2B event meanings. \\
Recursive self-inquiry source & Project-level question/method/implementation ratchet &
Adopted as the G3R bridge. It treats development as an RCI binding while keeping
candidate generation, independent return, review, and protected promotion separate. \\
Operator-first, regenerative-calculus, and compositional-grammar proposals &
Candidate relation/question kernels &
Archived as proposed derivations, not adopted as live authority. Their admissible
future role is bounded inert question generation from typed open relation ports after
measured failure evidence. Candidate completions are not truths; semantic completion is
not an external return; observation availability does not replace protected consequence;
and generic operator composition and semantic replacement remain deferred. \\
\bottomrule
\end{longtable}

The exact source bytes, original locations, archive paths, sizes, and SHA-256
digests are recorded in \path{docs/source-manifest.md}. No archived prompt or
document acquires authority merely by being stored in the repository.

\section{Aggregate, history, configuration, and retained state}

For a binding $B$, let $\mathsf E_B$ be the authoritative versioned event
vocabulary and let $\mathbf e_{\leq n}$ be an immutable ledger prefix. The
authoritative aggregate fold is
\[
 \Phi_{\mathrm{agg}}(\mathbf e_{\leq n})
 = \Sigma^{\mathrm{agg}}_{B,n},
 \qquad
 \Sigma^{\mathrm{agg}}_{B,n+1}
 = \operatorname{evolve}(\Sigma^{\mathrm{agg}}_{B,n},e_{n+1}).
\]
In the reference implementation, $\Sigma^{\mathrm{agg}}$ is
\texttt{InquiryState}. The fold is deterministic, effect-free under replay, and
reconstructs authoritative standing. It is not semantic compression.

A binding MAY define a partial derived relation
\[
 \rho_B:\Sigma^{\mathrm{agg}}_B\rightharpoonup\mathsf{Hist}_B
\]
from authoritative evidence to its realized interaction-history carrier. Ledger
succession and realized domain succession are distinct unless the binding proves
their coincidence for the records in question. No universal occurrence type is
assumed. A further binding-defined projection
\[
 p_B:\mathsf{Hist}_B\longrightarrow X_B
\]
maps realized history to a configuration carrier.

For a protected experiment horizon $\mathcal H$, exact history consequence
equivalence is
\[
 h\equiv_{\mathcal H}h'
 \quad\Longleftrightarrow\quad
 \forall e\in\mathcal H,\;
 \operatorname{Con}_e(h)=\operatorname{Con}_e(h'),
\]
where the binding supplies each consequence codomain and equality semantics. A
working discriminator relation $\approx_{\mathcal H,W}$ is weaker and MUST NOT
be promoted to exact equivalence merely because currently admitted probes have
not separated two histories.

A retained representation is a binding- and horizon-scoped map
\[
 q_{\mathcal H}:\mathsf{Hist}_B\longrightarrow S.
\]
It is consequence sufficient exactly when equality of represented states implies
protected history equivalence. Equivalently, every protected consequence map
factors through $q_{\mathcal H}$. Claiming the exact coarsest quotient additionally
requires the reverse implication and an independent minimality or isomorphism
argument. No finite size, computability, linearity, realizability, or bit
optimality follows merely from the quotient's set-theoretic existence.

The permanent architecture is
\[
 \boxed{\Phi_{\mathrm{agg}}\neq q_{\mathcal H}},\qquad
 \boxed{p_B\neq q_{\mathcal H}},\qquad
 \boxed{\tau_\pi\neq\mathsf{Hist}_B},\qquad
 \boxed{P\neq S_{\mathcal H}},
\]
where $\tau_\pi$ is one recurrent-probe trace and $P$ is a G2
\texttt{RetentionPackage}. These are semantic roles even if a binding happens to
reuse one underlying data type.

For every admitted continuation relation $\rho_a$ on history, a retained state is
recursively executable only when
\[
 q(h)=q(h')\Longrightarrow q[\rho_a(h)]=q[\rho_a(h')].
\]
Where extension by $u$ is deterministic, this can be checked as
\[
 q(h\mathbin{\cdot}u)=U(q(h),u).
\]
The reducer \texttt{evolve} and the retained-state update $U$ are distinct.
Consequence answer sufficiency and recursive state sufficiency are separate
claims.

Determination languages are carrier relative. A history determination descends
to retained state only when it is saturated under $q$, equivalently when it is
the inverse image of a determination over $S$. Frequent observation or a semantic
name does not establish descent.

The existing \texttt{PredictionSeal} is a sealed prospective consequence
commitment, not merely a scalar next-value guess. It MUST precede the actual
return used in comparison. The existing \texttt{Mismatch} remains a typed
comparison of a seal with an actual return and cannot select its own repair.
Updating $s$ under a fixed $q$, refining $q$, extending $\mathcal H$, and rebinding
$B$ are separate transitions. In particular, a G2B \texttt{MemoryPatchCandidate}
repairs a semantic lemma relation and is not a G3 representation refinement.

For a representation record $R=(q,\mathcal C,V,L,W)$, a strict successor
$R'\succ R$ requires preservation of every still-valid protected predecessor
competence, a typed strict gain, and standing independent warrant. An invalidated
predecessor competence requires an exact warranted disposition; a mere scope
tradeoff is incomparable rather than a strict successor. Candidate novelty,
newer model output, or regenerated skeptical prose is not gain. Active
representations are derived from immutable lineage and standing decisions; they
are not a new writable authority.

\section{Normative requirements}

\subsection{A. Constitution and semantic kernel}

\paragraph{RCI-001 --- Authority and traceability.}
The root specification MUST be the semantic authority. Every numbered
requirement MUST have one matrix disposition: implemented evidence, a named
later Goal, or explicit open research. Ambiguity MUST be recorded in the matrix
or an ADR and MUST NOT be silently resolved by code.

\paragraph{RCI-002 --- Typed binding and scope.}
Every inquiry MUST pin a binding revision, typed carriers and relations,
protected consequence profile and horizon, admissible operations, available
discharge mechanisms, assumptions, and local scope. Metrics, probability,
causality, time, utility, continuity, determinism, and observability MUST NOT be
assumed unless the binding declares them. Carrier role is part of semantic
interpretation. A binding that distinguishes configuration, realized history,
and retained representation MUST identify those roles in a versioned binding
manifest or later contract without mutating a sealed inquiry-start event schema.

\paragraph{RCI-003 --- Questions create claims.}
A question contract MUST bind referents, expose a typed relational role, delimit
lawful answer shape, and identify lawful follow-ups. Its answer MUST begin as a
provisional claim, never as a fact, formal candidate, external return, checker
verdict, or warrant.

\paragraph{RCI-004 --- Inert payload containment.}
Opaque L0 payloads, including empty, Unicode, contradictory, malformed-looking,
and prompt-injection strings, MUST remain representable as inert data. Payloads
and templates MUST NOT import code, execute tools, construct SQL or shell
commands, or select executable policy.

\paragraph{RCI-005 --- Lawful unknown and failure separation.}
\texttt{Unknown} is a lawful epistemic result. No attempt, unsupported
capability, timeout, transport failure, presentation uncertainty, capture
failure, malformed decode, logical unknown, checker indeterminacy, exhausted
search, and returned null MUST remain distinct. None proves impossibility,
necessity, sufficiency, equivalence, or control.

\paragraph{RCI-006 --- Description is not control.}
Descriptive factorization, correlation, prediction, witness of
may-reachability, or repeated observation MUST NOT create a control certificate.
Control requires an explicit action model, target, admissibility conditions, and
independently checked strategy or control evidence in the same scope.

\paragraph{RCI-007 --- Mandatory attack and non-explosion.}
Every proposed necessity, sufficiency, or prerequisite MUST create its
counterexample or alternate-route attack unless that exact proposition is
already discharged in the identical scope. Candidate contradiction MUST create
a deduplicated conflict obligation and MUST NOT permit explosion, arbitrary
inference, or arbitrary deletion.

\paragraph{RCI-008 --- Immutable succession.}
Claims, events, contracts, policies, bindings, lemma versions, predictions,
routes, returns, and retention packages MUST be immutable. Correction MUST
append typed relations or events such as \texttt{supersedes},
\texttt{splits\_from}, \texttt{merges\_from}, \texttt{rebinds},
\texttt{refutes}, \texttt{suspends}, \texttt{promotes}, or
\texttt{reopens}. Historical predecessors MUST remain addressable.
Ledger succession MUST NOT be interpreted as realized domain succession unless
the binding establishes their coincidence for the relevant records.

\paragraph{RCI-009 --- Capability and phase gates.}
The complete research catalog MAY be retained as versioned data, but only an
explicit milestone profile MAY be scheduled. A catalog entry, schema, command
group, optional dependency, or interface MUST NOT imply that its later-phase
semantics are implemented.

\paragraph{RCI-010 --- Consequence-conditioned relevance.}
Relative to a probe and protected profile, available structure MUST be
distinguished as active, relevance-undetermined, or warranted irrelevant.
\emph{Not active} MUST NOT mean \emph{irrelevant}. Aggressive occlusion requires
checked consequence-null standing and a reopening condition.

\subsection{B. Event spine, persistence, and ownership}

\paragraph{RCI-011 --- Pure transitions.}
The state machine MUST expose
\[
  decide(state,command)\rightarrow events,\qquad
  evolve(state,event)\rightarrow state.
\]
Both functions MUST be deterministic and free of I/O. Reducers MUST NOT
generate clocks, identifiers, randomness, provider metadata, or external work.

\paragraph{RCI-012 --- Durable authority and one owner.}
The append-only event ledger plus content-addressed artifacts is durable
authority. Within folded state, every semantic object MUST have one canonical
owner: episodic memory \(M_E\), semantic memory \(M_S\), procedural memory
\(M_P\), latent/retention memory \(M_L\), support/warrant structure \(W\),
attempt structure \(A\), or prediction structure \(\Pi\). Other occurrences
MUST be references or regenerable views, never independently writable truth.

\paragraph{RCI-013 --- Replay-complete records.}
Events and pinned manifests MUST retain every input needed for deterministic
replay, including binding revision, scope, assumption/guard AST, finite-universe
hash, explicit \texttt{closed\_world}, question/contract manifest and digest,
scheduler and warrant-policy versions, provenance, route snapshot, and artifact
references.

\paragraph{RCI-014 --- Rebuildable projections and snapshots.}
Active theory, candidate graphs, probe views, semantic fields, reopening views,
ancestry indexes, scheduler views, debt summaries, indexes, and reports MUST be
deterministic projections. Materialized projections MUST pin source sequence
and projection schema version. Projection rows and checkpoints MUST commit
together. Deleting projections and snapshots MUST NOT change replayed state or
deterministic export.

\paragraph{RCI-015 --- SQLite concurrency.}
The reference store MUST use stdlib SQLite in WAL mode, one writer path per
inquiry, \texttt{BEGIN IMMEDIATE}, expected stream sequence, and unique
event/entity identifiers. A racing append loser MUST receive a typed optimistic
concurrency conflict and MAY retry only after rebuilding from committed state.

\paragraph{RCI-016 --- Schema evolution.}
Every event MUST carry an explicit type and version. The greenfield database
begins at schema version 1, reserves a pure upcaster registry, and fails closed
on unknown versions. Upcasters MUST preserve meaning and MUST NOT fabricate raw
returns, closed support, or stronger warrants. Projections MAY be rebuilt;
authoritative events MUST NOT be rewritten.

\paragraph{RCI-017 --- Persisted effects and replay.}
External work MUST begin only from a committed effect-request event. Execution
MAY be at least once; acceptance MUST be idempotent. Replay MUST fold events
without executing an effect, consulting a network, generating new metadata, or
mutating an external system.

\paragraph{RCI-018 --- Request and attempt cardinality.}
A logical effect request and a delivery attempt MUST have different identities.
A request MAY have several attempts. Each attempt MUST have exactly one
terminal attempt outcome. A request MUST accept at most one resolution; late or
duplicate returns MAY be retained as rejected evidence but MUST NOT replace an
accepted resolution.

\paragraph{RCI-019 --- Plan and attempt outcomes.}
\texttt{NoAttempt} belongs to a plan disposition and MUST reference the plan,
not a nonexistent attempt. Actual attempt outcomes MUST distinguish at least
\texttt{NotPresented}, \texttt{PresentationUnknown}, \texttt{CaptureFailed},
\texttt{Cancelled}, and \texttt{Returned}, with typed timeout, transport, and
unsupported reasons where applicable.

\paragraph{RCI-020 --- Route identity.}
An allowlisted versioned \texttt{RouteDefinition} MUST be distinct from the
immutable per-attempt \texttt{RouteSnapshot}. The snapshot MUST record backend
and adapter versions, addressed source, redacted endpoint/channel identity,
execution-environment fingerprint, request/action digest, and ordered transform
manifest. Credentials MUST NOT be persisted. Addressed identity, source-reported
identity, and decoded semantic source claims MUST remain distinct.

\paragraph{RCI-021 --- Raw external return.}
An actual return MUST be captured before interpretation as immutable exact
bytes in content-addressed storage, with digest, size, media type, encoding, and
capture-boundary metadata. JSON null, empty bytes, empty string, zero, and false
MUST remain distinct. Artifacts MUST be published and hash-verified before an
event references them; orphan artifacts are recoverable, dangling references
are forbidden.

\paragraph{RCI-022 --- Decode, check, and warrant separation.}
An \texttt{ExternalReturn}, \texttt{DecodeOutcome}, \texttt{DecodedResult},
\texttt{CheckerVerdict}, \texttt{WarrantDecision}, and promotion MUST be
separate immutable records. Decode outcomes MUST include decoded, malformed,
unsupported, and failed. Checker outcomes MUST include valid, invalid,
indeterminate, timeout, unsupported, and failed. Adapter-reported success MUST
NOT be treated as independent checking.

\paragraph{RCI-023 --- Prediction and mismatch.}
A prediction used to assess mismatch MUST be sealed before its attempt can
produce a return and MUST pin target, scope, assumptions, allowed variation,
and basis. Only an actually returned and successfully decoded observation MAY
form semantic mismatch with a sealed prediction. Operational and decode
failures MUST NOT become fabricated mismatches. A prediction seal is a
prospective consequence commitment and MAY contain exact set-valued,
relational, bounded, or otherwise typed commitments; no scalar error is
required.

\paragraph{RCI-024 --- Memory and working views.}
\(M_E\) owns episodic records and artifact references, not copies of ledger
events; \(M_S\) owns immutable semantic lemma versions; \(M_P\) owns admitted
procedures, probes, contracts, and method manifests; \(M_L\) owns retention,
residue, and compression-application references. Active theory, probe state,
reopening state, working context, semantic field, and schedule MUST be derived.
The replay-complete aggregate \texttt{InquiryState} MUST NOT be treated as a
consequence-sufficient retained state. A probe trace is one comparable
observational subtrace, not a universal realized history.

\subsection{C. Claims, formalization, support, and warrant}

\paragraph{RCI-025 --- Orthogonal semantic axes.}
Claim lifecycle, assessment, representation level, execution outcome, decoded
logical result, checker verdict, warrant class, applicability, and active
promotion MUST be separate typed concepts. L3 MUST be a linked promoted view,
not an in-place mutation of an L0--L2 claim.

\paragraph{RCI-026 --- Structural L0 conflict.}
Opaque prose alone MUST NOT establish contradiction or entailment. L0
structural conflict is permitted only from explicit proposition identity,
opposite polarity, compatible role, identical bound referents, and identical
scope. Semantic contradiction requires validated reification or checked domain
evidence.

\paragraph{RCI-027 --- Restricted reification.}
Phase 2 formalization MUST use a serialized Boolean/finite-enum AST containing
only literals, symbol references, equality, negation, conjunction,
disjunction, implication, and equivalence. It MUST permit neither arbitrary
code nor quantifiers. Reification outcomes MUST be a strict union of reified,
needs clarification, unsupported, and failed. Reification alone carries no
warrant.

\paragraph{RCI-028 --- Support environment.}
A support environment MUST pin normalized assumptions, scope, binding revision,
finite universe where applicable, and an independent realizability verdict.
Absence of a known contradiction MUST NOT count as established realizability.

\paragraph{RCI-029 --- Route-specific dependency boundary.}
A \texttt{SupportRoute} MUST bind one conclusion and support environment to its
required dependencies, open dependencies, justifications, certificates,
checker, warrant decision, and provenance. Open dependency is defined per
route; it MUST NOT be flattened into one lemma-level list.

\paragraph{RCI-030 --- Minimal support and nogoods.}
Current minimal supports form subset antichains only among routes with the same
conclusion, scope, binding, applicability, universe, and policy versions.
Dominated routes MUST remain in history. A warranted nogood MUST deactivate
every environment containing it without deactivating unaffected alternate
routes.

\paragraph{RCI-031 --- Promotion policy.}
Hard promotion MUST be exact, scoped, applicable, independently checked,
dependency-closed, policy-authorized, and cycle-free. Model output, reification,
retrieval, unchecked backend assertions, and heuristic support MUST NOT hard
promote. A license or policy decision MUST reference warrant and MUST NOT be its
own warrant.

\paragraph{RCI-032 --- Support, proof, and ancestry topology.}
Ordinary active positive support MUST be grounded and acyclic; an ungrounded
self-supporting recurrence remains open. Mathematical cyclic proof material MAY
be stored as inert data, but only a later explicitly named independent cyclic
proof checker may turn it into a checked certificate leaf. Semantic succession
ancestry MUST use typed predecessor references and remain acyclic, including
split and merge histories.

\paragraph{RCI-033 --- Applicability and reopening.}
A lemma is active only when its applicability holds and at least one standing
minimal support route is closed. Guard, applicability, support, nogood, scope,
binding, policy, or retained-residue change MUST be able to deactivate and
reopen it. Deactivation MUST preserve prior history and provenance.

\paragraph{RCI-034 --- Lemma/version ownership.}
\(M_S\) MUST own the immutable semantic proposition, relation, scope,
applicability, and ancestry of a \texttt{LemmaVersion}; \(W\) MUST own its
support routes and warrant decisions. A \texttt{WarrantedLemmaView} MAY join
them but MUST NOT become a second authority.

\paragraph{RCI-035 --- Canonical result surface.}
The common decoded result surface MUST support \texttt{Witness},
\texttt{Counterexample}, \texttt{Separator},
\texttt{EquivalenceCertificate}, \texttt{Conflict},
\texttt{Prerequisite}, \texttt{ReachabilityWitness},
\texttt{UnreachabilityCertificate}, \texttt{Success}, \texttt{Failure}, and
\texttt{Unknown}. Each role MUST have a strict schema and canonical state
mapping. Unsupported later roles MUST fail closed.

\paragraph{RCI-036 --- Initial warrant classes.}
An independently evaluated witness MAY hard-warrant only the exact scoped
existential proposition it demonstrates. Exhaustive UNSAT MAY be hard only for
a declared closed finite universe. Z3-only UNSAT is
\texttt{solver\_trusted} and soft. Statistical or observational evidence MAY
support only the precisely qualified proposition its policy permits.

\subsection{D. Questions, scheduling, interfaces, and cognitive spine}

\paragraph{RCI-037 --- Versioned question catalog.}
The full question library inherited from the drafts MUST be retained as
immutable versioned catalog data. Contract lifecycle is
\texttt{draft -> experimental -> stable -> deprecated}; deprecated versions
remain replayable but MUST NOT be newly scheduled. Each inquiry MUST pin its
catalog manifest and digest.

\paragraph{RCI-038 --- Core profile.}
Only \texttt{core-v1} is schedulable in the Phase 1 kernel. It contains
obligation characterization, same-class variation, minimal boundary crossing,
factor proposal, necessity counterexample, sufficiency counterexample,
conflict localization, and residual characterization. Phase 2 MAY additionally
activate explicitly versioned warrant, localization, generalization,
actualization, and prerequisite profiles. Abstraction, recursive-formal,
compression, and control profiles remain inactive.

\paragraph{RCI-039 --- Deterministic scheduler.}
Obligation deduplication MUST key kind, deterministically normalized bound
arguments, scope fingerprint, and binding revision; the attempt key additionally
includes contract ID and version. Ready order MUST be invariant/warrant-conflict
safety, explicit priority, dependency depth, creation sequence, then stable ID.
Defaults are 100 reducer steps, three attempts per identical attempt key, and a
60-second external-effect timeout. Exhaustion yields unknown under the present
binding.

\paragraph{RCI-040 --- Public lifecycle.}
The SDK and CLI MUST provide semantically equivalent \texttt{start},
\texttt{step}, \texttt{run}, \texttt{resume}, \texttt{inspect},
\texttt{submit\_answer}, \texttt{replay}, and \texttt{export} inquiry flows,
plus contract, evaluation, database, and backlog command groups. Replay and
export MUST be deterministic and effect-free.

\paragraph{RCI-041 --- Generators and optional OpenAI adapter.}
Manual and scripted generators MUST support all blocking tests without a
network. The optional provider-neutral OpenAI Responses API adapter MUST use an
explicitly configured model, locally compiled stateless context, no model
tools, bounded output, \texttt{store=False}, preserved raw L0 answers, and
recorded execution status and usage. It MUST NOT be required for replay or
reference demonstrations.

\paragraph{RCI-042 --- Recurrent probes.}
Recurrent-probe identity MUST pin relational role, binding schema,
applicability, comparison semantics, and protected horizon. Same wording is not
sufficient; a changed binding requires an explicit checked comparison bridge.
Ordered comparable probe events form a queryable trace, but perceived
persistence, change, recurrence, covariation, or path dependence remains a
provisional claim.

\paragraph{RCI-043 --- Semantic field and fresh observation.}
The question-conditioned semantic field MUST be a derived generator view, never
an authority. A model relevance judgment MUST NOT suppress inquiry. Where an
environment can be independently observed, a fresh return SHOULD be captured
before prior semantic answers are exposed; comparison MAY follow after capture.

\paragraph{RCI-044 --- Cognitive learning boundary.}
The Phase 1--2 cognitive slice MUST separate episodic retention, decoded
interpretation, candidate reconstruction, semantic delta, and warranted
semantic promotion. Reconstruction detail is neither history nor knowledge.
Full relational retrieval, consolidation, reconsolidation, retention economics,
and learned-probe promotion are Phase 2-following work.

\paragraph{RCI-045 --- Circuit reference binding.}
The deterministic eight-state circuit
\[
 lamp\_on = switch\_closed \land (main\_power \lor backup\_power)
\]
MUST show a backup-power counterexample to main-power necessity, an open-switch
counterexample to available-power sufficiency, and exhaustive closed-world
establishment of switch-closed necessity.

\paragraph{RCI-046 --- Route-graph reference binding.}
The finite transition binding MUST contain
\texttt{start -> gate -> target} and
\texttt{start -> bypass -> \{target, dead\_end\}}. It MUST show a route
bypassing the proposed gate prerequisite and distinguish may-reachability from
must-control; nondeterministic bypass never promotes to must-control.

\subsection{E. Consequence quotient, retention, and compression}

\paragraph{RCI-047 --- Permanent no-collapse distinctions.}
The core semantics MUST distinguish exact consequence quotient from lossy
approximation, source/object reconstruction error from protected consequence
loss, and present use from reconstruction, direct consequence evaluation, and
reacquisition. Reconstruction uses retained state and admissible computation
without substantial new external learning; reacquisition uses new evidence,
instruction, interaction, or practice. None of those routes implies another.
Exactness is always relative to a pinned protected horizon and scope, never
global by default.
Configuration, realized-history, and retained-state carrier roles MUST remain
distinct even when a binding reuses an underlying data type. Configuration
projection and history quotient are different contracts. A G2
\texttt{RetentionPackage} and a G2B \texttt{MemoryPatchCandidate} MUST NOT be
treated as a certified retained state or a representation refinement.

\paragraph{RCI-048 --- Compression decision layers.}
\texttt{CompressionContract}, \texttt{CompressionValidation},
\texttt{CompressionLicense}, and \texttt{CompressionApplication} MUST be
separate immutable concepts. Exact and approximate licenses MUST be a
discriminated union, not a Boolean plus optional-field bag. A measured loss of
zero MUST NOT be relabeled exact without universal or closed-domain separating
evidence. Every contract MUST pin explicit source and target carrier roles and
schemas, binding, scope, protected horizon, continuation family, consequence
family, equality semantics, and recovery semantics. Exact history-state
validation MUST separately check consequence factorization and every claimed
continuation descent.

\paragraph{RCI-049 --- Retention capability routes.}
A retention package MUST reference its representation artifacts and compression
application when one exists. It MAY independently expose a
\texttt{DirectUseRoute}, \texttt{ObjectRegenerationRoute},
\texttt{ConsequenceEvaluationRoute}, or \texttt{ReacquisitionRoute}. Each route
MUST pin its \texttt{RecoveryLicense}, fidelity or loss semantics, protected
query/competence profile, implementation version, support route, dependencies,
provenance, failure behavior, and fallback. A reacquisition route MUST reference
a \texttt{ReacquisitionScaffold} rather than claim to contain the recovered
competence. Combined modes are derived from explicit routes, not stored as an
ambiguous enum or optional-field bag.
A G2 \texttt{RetentionPackage} is only a typed package of references and
provisional routes. Any route consuming a compressed representation MUST
identify the independently validated compression application and standing
recovery license that make it an exposed capability.

\paragraph{RCI-050 --- Residue, debt, and reopening.}
Compression MUST return retained representation plus dependency residue,
regeneration/evaluation routes, provenance, authoritative ancestry, debt,
fallback, and reopening predicates. It MUST NOT erase an open dependency or
assume that a disappeared representation dependency was discharged.
Approximate bounds compose only through a pinned checked composition policy;
otherwise composition fails closed.
Representation/path residue and epistemic open support dependency are distinct
records and MUST NOT discharge or substitute for one another.

\paragraph{RCI-051 --- Exact linear family theorem.}
Let \(Q\subseteq\mathbb R^d\), \(V=\operatorname{span}Q\), and
\[
 x\sim_Qy \iff q^\top x=q^\top y\quad\text{for every }q\in Q.
\]
Then
\[
 x\sim_Qy \iff x-y\in V^\perp,\qquad
 \mathbb R^d/V^\perp\cong V
\]
through \([x]\mapsto P_Vx\). This is an exact worked binding, not a universal
Euclidean axiom.

\paragraph{RCI-052 --- Distributional and vector-output corollaries.}
For random \(q\) with \(\mathbb E\|q\|^2<\infty\), let
\(M=\mathbb E[qq^\top]\). Then
\[
 q^\top x=q^\top y\ \text{almost surely}
 \iff x-y\in\ker M,
\quad
 \mathbb R^d/\ker M\cong\operatorname{im}M.
\]
This is almost-sure, not universal, equivalence unless the support conditions
justify the latter. For measurable vector-output \(A_q\), require
\(\mathbb E\|A_q\|_F^2<\infty\) and use
\(G=\mathbb E[A_q^\top A_q]\) with
\(\mathbb E\|A_q(x-y)\|_2^2=(x-y)^\top G(x-y)\).

\paragraph{RCI-053 --- Finite probes and minimality.}
For stacked finite probes \(A\) and diagonal \(W\succ0\),
\[
 M=A^\top W A,\qquad \ker M=\ker A,\qquad
 \operatorname{rank}M=\operatorname{rank}A.
\]
A new scalar probe adds one exact dimension exactly when its row lies outside
the current row span. \(\operatorname{rank}M\) is the minimum dimension among
linear encoders preserving all protected linear consequences; no unrestricted
bit-complexity or nonlinear minimality claim follows.

\paragraph{RCI-054 --- Correct reopening condition.}
A previously invisible direction becomes visible exactly when some
\[
 v\in\ker M_t\quad\text{but}\quad v\notin\ker M_{t+1},
\]
equivalently
\[
 \boxed{\ker M_t\not\subseteq\ker M_{t+1}}.
\]
When only positive-weight probes are added,
\(\ker M_{t+1}\subsetneq\ker M_t\). If the discarded distinction has no
residue or reacquisition route, reopening MUST return Unknown rather than
fabricating regeneration.
Generically, reopening occurs whenever a new protected consequence fails to
factor through the incumbent representation, equivalently when two inputs
identified by the incumbent have different new consequences. This criterion
applies independently of vector or linear structure.

\paragraph{RCI-055 --- Approximate geometry is not exact proof.}
For reconstruction error \(e=x-\hat x\),
\(e^\top M e=\sum_i\lambda_i(U^\top e)_i^2\) describes consequence
sensitivity under the declared quadratic loss. It does not solve optimal bit
allocation, which also depends on source/query distributions, code class,
decoder, randomness, and resource model. Numerical near-zero eigenvalues,
empirical zero loss, unbiasedness in expectation, and average bounds MUST NOT
establish exact nullity or worst-case safety.

\paragraph{RCI-056 --- Native compression adapters.}
Task-based and indirect rate-distortion theory MAY supply native approximate
bindings. EDEN, DRIVE, TurboQuant, and QJL are method-catalog candidates, not
RCI primitives. Native execution is later and MUST use an allowlisted,
digest-pinned, nonroot, network-disabled, read-only, capability-dropped,
resource-bounded container or equivalently isolated adapter. It MUST receive no
credentials, Docker socket, ledger/CAS authority, or live source mount.
Missing runtime or GPU is unsupported/no-attempt, not epistemic Unknown.

\subsection{F. Milestones, governance, and verification}

\paragraph{RCI-057 --- G1 boundary.}
G1 comprises Foundation, specification Phases 1--2, and the bounded cognitive
spine: packaging/CI/CLI; immutable events and pure reducers; SQLite/CAS,
snapshots, projections, and replay; question/claim/conflict/obligation kernel;
manual/scripted generators; restricted AST and exhaustive checking; optional
OpenAI/Z3; support routes and warrant; recurrent probes, prediction, raw
returns, reconstruction, semantic delta; both reference bindings; and
shadow-only backlog dogfooding.

\paragraph{RCI-058 --- G2 boundary.}
G2 adds relational retrieval, reconstruction policy, consolidation and
reconsolidation, retention without compression, semantic-field evaluation,
candidate learned probes, reacquisition-scaffold activation, and reacquisition
as an ordinary inquiry. It establishes pinned baseline and multidimensional
recovery-cost observations without claiming a universal scalar score. Learned
automata, probes, and representations remain hypotheses until separately
challenged and warranted.

\paragraph{RCI-059 --- G3 boundary.}
G3A begins with explicit source/target carrier roles, history consequence
equivalence, exact factorization, continuation compatibility, deterministic
recursive update where applicable, determination descent, path residue, generic
reopening, route-specific \texttt{RecoveryLicense} enforcement, and
representation succession. Its history-state slice precedes the exact
rational finite/linear bindings. G3B implements approximate loss licenses,
licensed recovery-frontier budgets, risk/confidence/debt composition, numerical
candidates, and a deterministic toy quantizer/oracle. G3C adds isolated native
adapters after license and implementation provenance are verified.

\paragraph{RCI-060 --- Later research boundaries.}
G4 covers recursive formal machinery such as CHC and PDR; G5 covers controller
synthesis and checked feedback strategies; G6 covers multi-backend evidence;
G7 covers compatibility migrations, hardening, security, performance,
evaluation, the opaque controlled-memory end-to-end benchmark, stable SDK, and
release automation. None is G1 acceptance.

\paragraph{RCI-061 --- Human-owned backlog policy.}
\path{.rci/config.toml} is tracked human-owned authority; runtime database,
artifacts, projections, and exports are local ignored state. Backlog reconcile
is shadow/dry-run by default. G1 manual apply MAY append only allowlisted
create, exact-dedupe, rank, and block effects. Close is proposal-only until a
later explicit human policy and Goal decision; tests cannot grant authority.
A regression appends a linked recurrence rather than rewriting closed history.

\paragraph{RCI-062 --- Evidence-runner authority.}
Evidence runners MUST use explicit argument vectors, bounded time and output, a
captured temporary workspace, and no network by default. They MUST expose no
source-writing, Git mutation, merge, push, policy editing, packaging,
deployment, release, or authority-expansion capability.

\paragraph{RCI-063 --- Blocking evidence.}
G1 MUST test arbitrary payload containment; reducer illegal transitions;
effect-free byte-identical replay/export; artifact tamper and crash points;
SQLite optimistic races; request/attempt cardinality, timeout, duplicates, and
first acceptance; null/empty/false payload distinctions; decode/check/warrant
separation; environment realizability; support and ancestry cycles; guard
deactivation; probe comparability and fresh-observation isolation;
prediction-before-return; reconstruction/history/knowledge separation;
unknown-relevance preservation; the two reference conclusions; SDK/CLI parity;
and backlog non-mutation and authority limits.

\paragraph{RCI-064 --- Honest conformance.}
Verification commands and exact returns MUST be recorded in
\path{docs/verification.md}. A failed, skipped, unavailable, or unrun check MUST
remain visible. No milestone may close because code looks plausible, because a
stub returns a passing shape, or because a later capability has a catalog
entry. Missing credentials and unavailable optional services do not block an
offline milestone; affected live claims remain unverified.

\subsection{G. Recovery and future benchmark addendum}

\paragraph{RCI-065 --- Recovery modes are relational.}
For retained state \(m\), context \(c\), protected horizon \(\mathcal H\), and
target competence \(k\), the system MUST represent present use,
reconstruction, direct consequence evaluation, and reacquisition as independent
route relations rather than nominal memory-object kinds. A recovered competence
\(k'\) need only be consequence-equivalent to \(k\) over the pinned protected
horizon; historical representation identity is not required. Failure of recall
or reconstruction MUST NOT be interpreted as absence of retained learning. A
route over retained material MAY consume a licensed history-state
representation, but it acquires no such license from the retained package alone.

\paragraph{RCI-066 --- Reacquisition advantage is baseline-relative.}
For each evaluated competence, context, and horizon, the baseline state or
cohort, protocol, evaluator, cost dimensions, budget, and comparison policy MUST
be pinned. The attainable recovery-cost frontier
\(\mathfrak C_{\mathcal H}(m,c,k)\) MAY include latency, interaction, compute,
external evidence, risk, and other binding-declared costs. Reacquisition
advantage requires a checked strict Pareto improvement, or a binding-authorized
ordering, relative to the pinned baseline frontier
\(\mathfrak C^0_{\mathcal H}(c,k)\). No universal scalarization is assumed; a
scalar savings score is lawful only when the binding explicitly supplies one.
Competence equivalence MUST use a binding-declared typed relation. History
consequence equivalence may be reused only when the competence values inhabit
that declared history carrier.

\paragraph{RCI-067 --- Recovery licence, scaffold, and forgetting.}
A \texttt{RecoveryLicense} MUST pin consequence fidelity, latency, computation,
external evidence, interaction, risk, resources, support, scope, policy version,
and failure semantics required by its route. A
\texttt{ReacquisitionScaffold} MAY retain cues, probes, representations,
methods, boundaries, failures, prerequisites, and provenance, but it counts as
retained learning only when checked evidence shows a licensed recovery
advantage. Forgetting is a reduction in future recovery capacity, not merely
deletion. Reopening MAY activate present use, reconstruction, provenance
retrieval, or a new reacquisition inquiry; it MUST NOT fabricate a route that
was not retained and warranted.

\paragraph{RCI-068 --- Opaque controlled-memory benchmark boundary.}
The opaque controlled-memory environment is a future executable G7 benchmark,
not a core data type, adapter requirement, or G1 blocking test. Its staged
capabilities include authoritative raw returns, intervention/return alignment,
temporal and derived probes, sealed prediction, conditional/support-route
discovery, predictive representation, control, consequence-relative
compression, reopening, relearning advantage, and transfer. Earlier Goals MAY
deliver individually scoped prerequisites without claiming the end-to-end
benchmark. Predictive State Representation is a future native binding rather
than an RCI reinvention. The exact linear consequence theorem applies only after
a warranted vector representation and linear protected query family have been
established; it MUST NOT be imposed on raw memory or an arbitrary nonlinear
environment.
The benchmark MUST include two realized histories with the same configuration
projection but different protected future consequences, and it MUST demonstrate
discovery and prospective use of the missing path distinction before claiming
learned predictive state.

\paragraph{RCI-069 --- Aggregate and retained-state folds do not collapse.}
The project MUST distinguish the authoritative replay fold
$\Phi_{\mathrm{agg}}:\mathsf E^*\to\Sigma^{\mathrm{agg}}$ from a
binding-defined retained representation
$q_{\mathcal H}:\mathsf{Hist}_B\to S_{\mathcal H}$. \texttt{InquiryState},
\texttt{ProbeTrace}, and \texttt{RetentionPackage} MUST NOT be treated as
certified minimal consequence-sufficient state merely because they fold or
retain prior structure. The ledger/CAS remains replay-complete authority even
when a licensed semantic representation removes protected-irrelevant history
distinctions.

\paragraph{RCI-070 --- Recursive state sufficiency.}
A representation claimed as executable retained state under an admitted
continuation family MUST establish that equal represented states have identical
sets of represented successors for every admitted operation. Where extension is
deterministic, an independently checked update law MAY establish this as
$q(h\mathbin{\cdot}u)=U(q(h),u)$. Failure MUST reject the recursive-state claim
or narrow its continuation scope without weakening a separately valid answer-
sufficiency claim.

\paragraph{RCI-071 --- Representation successor ratchet.}
A generated or newer representation MUST NOT replace an incumbent by novelty or
recency. Strict succession requires preservation of every still-valid protected
predecessor competence, at least one typed strict gain, and standing independent
warrant. Independently invalidated predecessor claims require exact warranted
dispositions. Scope tradeoffs and candidates that dominate in neither direction
remain explicit frontier members or provisional candidates; no total order is
assumed.

\section{Recursive project inquiry}

The repository and its operating agent MAY be investigated through an ordinary
project binding. This does not make model weights, platform policy, credentials,
or external account authority writable RCI state. It makes the development argument
replayable:
\[
 \text{anchor}\to\text{limitation}\to\text{candidate frontier}
 \to\text{sealed Goal}\to\text{external return}\to\text{review}
 \to\text{successor decision}\to\text{protected promotion}.
\]
Changes to theory, question/probe repertoire, method repertoire, representation,
Goal decomposition, and implementation are different successor kinds. A failure in
one kind does not authorize deformation of another kind.

\paragraph{RCI-072 --- Project binding and exact anchor.}
Every recursive development cycle MUST begin from a clean exact repository commit,
tree digest, live-authority digest, gate digest, protected branch, and repository
identity. Stale or dirty anchors fail closed. The project binding reuses the event
ledger/CAS and MUST NOT create a second authoritative development history.

\paragraph{RCI-073 --- Consequential capability limitation.}
A project limitation MUST name current and missing capability, its typed limitation
kind, preserved observation evidence, protected consequences, and the boundary across
which those consequences differ. Mere novelty, dissatisfaction, generated skepticism,
or roadmap position is not a limitation. If no possible return changes a protected
successor state, the question is inert.

\paragraph{RCI-074 --- Question-repertoire succession.}
A generated question contract is an inert candidate. Admission requires typed
referents, explicit preconditions, at least two consequentially distinct possible
returns, comparison semantics, downstream consumers, falsifying attacks, versioned
controller policy, and independent evidence. Initially admitted project questions are
confined to the \texttt{recursive-project-v1} binding profile and cannot suppress
ordinary inquiry or execute arbitrary code.

\paragraph{RCI-075 --- Method-repertoire succession.}
A candidate method MUST bind a surviving relation to the established field in which
it is native, preserve primary-source artifacts and licence, state assumptions and
applicability checks, and report whether an adapter is missing. Admission does not
install a dependency or implementation. Missing executable support opens a separate
sealed implementation Goal.

\paragraph{RCI-076 --- Goal and implementation succession.}
An implementation Goal MUST be sealed before source change and identify
\(Current,Desired,Separator,Preserve,Acceptance,Scope,Assumptions\), expected incumbent
and candidate returns, allowed mutation roots, forbidden authority roots, incumbent
and proposed gate digests, rollback, and reopening. A project successor replaces an
incumbent only if it preserves or explicitly dispositions predecessor capabilities,
demonstrates a typed strict gain, and has exact independent evidence; otherwise it is
kept, rejected, or retained on an incomparable frontier.

\paragraph{RCI-077 --- Candidate actualization and protected promotion.}
Candidate development occurs from the exact anchor in an isolated branch/worktree.
The candidate cannot change its sealed Goal, erase failing evidence, review itself,
or promote itself. Fresh review and hosted CI bind the exact candidate SHA. The runtime
MAY record externally performed promotion but MUST expose no source-writing, arbitrary
command, Git mutation, credential, force-push, release, deployment, or permission-
expansion port. Required-check evolution uses a dual-gate overlap before replacement.

\paragraph{RCI-078 --- Bounded continuous recurrence.}
After a verified cycle, retain its evidence, regression boundaries, failures,
reacquisition scaffold, cleanup, and consequential residue, then open at most one next
bounded Goal. Stop or return \texttt{Unknown} when there is no consequential residue,
no lawful discriminator, the same blocker recurs three times, evidence is invalid or
indeterminate, or progress requires authority outside the active Goal. Recurrence MUST
NOT relax its own stopping or acceptance rules merely to continue.

\paragraph{RCI-079 --- Bounded review equivalence.}
A model-disconnected review route MAY establish only
\texttt{valid\_within\_profile} relative to a closed, versioned fault family and exact
Goal, candidate-environment, candidate-commit, gate, invariant, evidence, and reproducer
pins. Every required seeded fault MUST be observed through an independent exact-head
development return. A surviving seeded fault makes the assessment invalid; missing,
stale, foreign, malformed, unsupported, or indeterminate evidence makes it
indeterminate. Passing the profile leaves unbounded semantic review \texttt{Unknown}.
The assessment MUST NOT be relabelled as an \texttt{IndependentReview}, satisfy a
fresh-review requirement, warrant a project successor, or authorize promotion.
Model-generated semantic breakers are strictly parsed inert candidates; malformed model
output remains indeterminate and cannot weaken the deterministic profile. This
requirement is accepted semantically but remains unverified after the stopped G3V
candidate.

\section{Consequence operators retained from prior drafts}

For a typed relation \(R:X\rightsquigarrow Y\) and target predicate
\(P\subseteq Y\), RCI uses
\[
 \Diamond_R(P)=\{x:R[x]\cap P\neq\varnothing\},\qquad
 \Box_R(P)=\{x:R[x]\subseteq P\},
\]
and, when execution must exist,
\[
 \operatorname{Must}_R(P)=\operatorname{Dom}(R)\cap\Box_R(P).
\]
For \(R:X\rightsquigarrow Y\) and \(S:Y\rightsquigarrow Z\), composition
retains the order
\[
 \Diamond_{S\circ R}=\Diamond_R\circ\Diamond_S,\qquad
 \Box_{S\circ R}=\Box_R\circ\Box_S.
\]
These equations are normative and MUST be protected by tests.

\section{Compressed governing laws}

\begin{enumerate}[label=\arabic*.]
\item Questions create typed claims, not facts.
\item Repeated comparable probes create provisional perception, not warrant.
\item External return is not interpretation, and interpretation is not warrant.
\item Every \emph{must}, \emph{enough}, or \emph{prerequisite} opens its
      falsifying attack.
\item Contradiction creates inquiry, not explosion.
\item Description and may-reachability do not create control.
\item Hard knowledge requires closed, exposed, independently checked support.
\item Active representation may revise; authoritative event history does not.
\item Compress only under an exact quotient or explicit approximate license,
      while preserving residue and reopening.
\item Treat use, reconstruction, consequence evaluation, and reacquisition as
      separate licensed recovery paths; forgetting is reduced recovery capacity.
\item Every consequential residual becomes the next bounded obligation.
\end{enumerate}

\section*{Normative companion records}

The following repository records are normative at their respective authority
levels and MUST stay coherent with this specification:
\begin{itemize}
\item \path{PLAN.md}: approved engineering architecture and delivery sequence;
\item \path{docs/architecture.md}: implementer-facing ownership and data flow;
\item \path{docs/requirements-matrix.md}: requirement disposition and evidence;
\item \path{docs/adr/}: explicit reconciliations and consequential decisions;
\item \path{docs/goals/}: bounded milestone targets, including sealed G1/G2A/G2B
      and the active G3A history-state target;
\item \path{docs/verification.md}: executed checks and honest returns;
\item \path{docs/source-manifest.md}: immutable source provenance.
\end{itemize}

\end{document}
